\documentclass[12pt]{article} 
\renewcommand{\thempfootnote}{\arabic{mpfootnote}}

\usepackage[hidelinks]{hyperref}

\usepackage[utf8]{inputenc}

\usepackage[T1]{fontenc}
\usepackage[ttdefault=true]{AnonymousPro}

\usepackage{multicol}
\setlength{\columnsep}{1cm}

\usepackage{abc}
\usepackage{graphicx}
\usepackage{musicography}

\usepackage[margin=0.5in]{geometry} 
\geometry{letterpaper}

\usepackage{fancyhdr} 
\pagestyle{fancy} % options: empty , plain , fancy
\renewcommand{\headrulewidth}{0pt} % customise the layout...
\lhead{}\chead{}\rhead{}
\lfoot{}\cfoot{\thepage}\rfoot{}

\begin{document}

\thispagestyle{empty}

\begin{abc}[name=speedtheplough-berry,options={-O= +c -F single.fmt}]
X:1
T:Speed the Plough
C:from Glenn Berry, Seabeck, WA
M:C|
K:D
|: AB/c/ | "D"d2dc dAF2 | dFcF BFAF | "G"GAGD "D"FGFD | "A7"E2EF EFED | 
"D"d2dc dAFA | dFcF BFAF | "G"GAGD "D"FGFD | "A7"EDCE "D"D2 :|
|: D2- | "D"DEFG ABAF | ABAF AdAF | "G"GBGA "D"FGFD | "A7"E2EF EFEE | 
"D"DEFG ABAF | ABAF AdAF | "G"GBGA "D"FGFD | "A7"EDCE "D"D2 :|
\end{abc}
\vskip 0.5in
\noindent\makebox[\textwidth][c]{%
  \begin{minipage}{7.0in}
    \setlength{\parindent}{1.5em}

    As with many other old tunes, Speed the Plough is so old and well
    established, no one is quite sure where (or when) it came
    from. This tune is associated with northern England and ``Plough
    Monday'', a holiday dating back to the late 15th century marking
    the first Monday after Epiphany, January 6th, that signals the
    beginning of the English agricultural year.\footnote{
      \url{https://tunearch.org/wiki/Annotation:Speed\_the\_Plough\_(1)}
      and \url{https://en.wikipedia.org/wiki/Plough\_Monday} \\
      Plough Monday is still celebrated in Whittlesea, Peterborough, England
      with the appearance of Whittlesea Straw Bear and the Whittlesea
      Straw Bear Festival (Instagram as @whittleseastrawbear ).} This
    tune is said to have been composed in 1799 by John Moorehead, born
    in Edinburgh and associated with county Armagh, Ireland. Shortly
    afterwards the tune is linked to the play, \textbf{Speed the
      Plough}, by Thomas Morton (1800).\footnote{tunearch.org,
      \textit{loc. cit.}} Graham Dixon's 1995 book, \textbf{Remember
      Me, The Fiddle Music of Robert Whinham}, includes a variation
    set with this tune, Robert Whinham being a dancing master and
    fiddler from Morpeth, Northumberland (1814 –
    1893).\footnote{Ibid. and
      \url{https://www.abebooks.com/servlet/BookDetailsPL?bi=22852432162&cm_sp=rec-_-pd_hw_i_1-_-bdp&reftag=pd_hw_i_1}}
    Meanwhile in America, the tune is found in George E Blake’s
    \textbf{Blake’s New and Complete Preceptor for the Violin}
    (Philadelphia, 1810?), George P Knauff’s \textbf{Virginia Reels}
    (Baltimore, 1839), and Elias Howe Jr’s \textbf{School for the
      Fife} (Boston, 1843).\footnote{Goertzen, Chris. \textbf{George P
        Knauff’s Virginia Reels and the History of American Fiddling},
      p 35, p 45.} Speed the Plough is still played on both sides of
    the Atlantic, being popular in New England (contradance), the
    Ozark region (old time), and throughout America and Canada in
    Irish traditional music sessions.\footnote{Ibid., 47, and
      \url{https://thesession.org/tunes/1191}}
    
  \end{minipage}
}


\end{document}
